\documentclass[aspectratio=169]{beamer}
% Theme and colors
\usetheme{SimplePlus}
\usecolortheme{default}

% Useful packages
\usepackage[utf8]{inputenc}
\usepackage{graphicx}
\usepackage{amsmath}
\usepackage{amssymb}
\usepackage{listings}
\usepackage{xcolor}

\newcommand{\quotes}[1]{``#1''}

% Monospace style for short code snippets
\lstset{
  basicstyle=\ttfamily\footnotesize,
  breaklines=true,
  columns=fullflexible,
  backgroundcolor=\color{black!5},
  frame=single,
  rulecolor=\color{black!20},
  xleftmargin=4pt,
  xrightmargin=4pt,
  aboveskip=4pt,
  belowskip=4pt
}

% Presentation information
\title{Smart Ambulance IoT}
\subtitle{Internet of Things}
\author{Giuseppe Di Blasi}
\institute{Università di Parma}
\date{}

\begin{document}

% Title slide
\begin{frame}
\titlepage
\end{frame}


% ======================================================
\section{Introduction}
% ======================================================

% \begin{frame}{The Smart Ambulance Scenario}
% \begin{itemize}
% \item An ambulance $\mathcal{A}$ with two operating environments: \textbf{Cockpit} ($\mathcal{A}_C$) and \textbf{Rear Cabin} ($\mathcal{A}_R$)
% \item Two first-aid workers acting as CoAP clients: \textbf{Driver} and \textbf{Reliever}
% \item A \textbf{GPS} receiver, also modeled as a CoAP client
% \item Each environment is a CoAP server hosting three resources: temperature sensor, air-conditioning controller, HMI
% \end{itemize}
% \begin{alertblock}{Focus of this talk}
% Not the CoAP theory itself, but the concrete \textbf{implementation choices} made to satisfy the project requirements.
% \end{alertblock}
% \end{frame}

% ======================================================
\section{Architecture}
% ======================================================

\begin{frame}{Two CoAP Servers, Three CoAP Clients}
\begin{columns}
\column{0.5\textwidth}
\textbf{Servers}
\begin{itemize}
\item \texttt{Cockpit} $\rightarrow$ port 5684
\item \texttt{RearCabin} $\rightarrow$ port 5683
\item Both host the same three resources: \texttt{TemperatureSensor}, \texttt{AirConditioningController}, \texttt{HMI}
\end{itemize}
\column{0.5\textwidth}
\textbf{Clients}
\begin{itemize}
\item \texttt{Driver}, \texttt{Reliever}, \texttt{GPS}
\item All extend a common abstract class \texttt{SmartAmbulanceClient}
\item Shared reusable GET/PUT operations toward HMI, sensor and controller
\end{itemize}
\end{columns}
\vspace{0.4cm}
\begin{block}{Design choice}
Two independent \texttt{CoapServer} instances instead of one server with parametrized resources — keeps the two environments fully separated.
\end{block}
\end{frame}

% \begin{frame}{Resources as POJO Classes}
% \begin{itemize}
% \item Each resource extends \texttt{CoapResource} and keeps its state in \textbf{private fields}
% \item State is exposed only via \texttt{handleGET} / \texttt{handlePUT} overrides — as requested in the project's ``General Suggestions''
% \item No module ever touches another class's internal variables directly
% \end{itemize}
% \vspace{0.3cm}
% \begin{exampleblock}{Resources implemented}
% \begin{itemize}
% \item \texttt{TemperatureSensor} — GET current temperature; observable
% \item \texttt{AirConditioningController} — GET target temp, PUT with permission check
% \item \texttt{HMI} — GET emergency state, PUT with cross-environment propagation
% \end{itemize}
% \end{exampleblock}
% \end{frame}

% ======================================================
\section{Implementation Choices}
% ======================================================

\begin{frame}[fragile]{Message Format: Hand-Built JSON}
\begin{columns}
\column{0.48\textwidth}
\begin{itemize}
\item Payloads built with \texttt{String.format}
\item Parsed with a small helper: \texttt{extractJsonValue(json, key)}
\item Every payload carries a \texttt{"Client"} field identifying the sender
\end{itemize}
\column{0.52\textwidth}
\begin{lstlisting}
{
  "Client": "DRIVER",
  "temp_target": 21
}
\end{lstlisting}
\begin{lstlisting}
extractJsonValue(
  payload, "temp_target")
\end{lstlisting}
\end{columns}
\begin{block}{Why it matters}
The \texttt{Client} field is what the server uses to route authorization logic (next slide).
\end{block}
\end{frame}

\begin{frame}[fragile]{Authorization Rules on the AC Controller}
\begin{columns}
\column{0.5\textwidth}
\textbf{Driver}
\begin{itemize}
\item Cockpit $\rightarrow$ \textcolor{green!50!black}{allowed}
\item Rear Cabin $\rightarrow$ \textcolor{green!50!black}{allowed}
\end{itemize}
\column{0.5\textwidth}
\textbf{Reliever}
\begin{itemize}
\item Cockpit $\rightarrow$ \textcolor{red}{FORBIDDEN}
\item Rear Cabin $\rightarrow$ \textcolor{green!50!black}{allowed}
\end{itemize}
\end{columns}
\vspace{0.3cm}
\begin{lstlisting}
if (client.equals("RELIEVER") && port == COCKPIT_PORT) {
    exchange.respond(ResponseCode.FORBIDDEN);
    return;
}
\end{lstlisting}
\begin{alertblock}{Note}
The check is done inside \texttt{handlePUT}, comparing the \texttt{Client} field with the server's own port.
\end{alertblock}
\end{frame}

\begin{frame}{Dynamic Temperature + Observe Pattern}
\begin{columns}
\column{0.5\textwidth}
\textbf{Gradual convergence}
\begin{itemize}
\item \texttt{TimerTask} with \texttt{scheduleAtFixedRate}
\item Current temperature moves $\pm1^\circ$C every 1.5s toward the target
\item Triggered by \texttt{setTargetTemperature()}
\end{itemize}
\column{0.5\textwidth}
\textbf{CoAP Observe}
\begin{itemize}
\item A client can observe the temperature sensor
\item The subscribed client is updated at every change
\item Unsubscribes with \texttt{stopObservingTemperature()}
\end{itemize}
\end{columns}
\vspace{0.4cm}
\begin{block}{Design choice}
Simulating a realistic air-conditioning response instead of an instantaneous jump to the target temperature.
\end{block}
\end{frame}

\begin{frame}[fragile]{HMI Synchronization Between Environments}
\begin{itemize}
\item When Driver or Reliever update emergency level / location,  the \texttt{HMI} resource itself opens an internal \texttt{CoapClient}
\item Same thing applies for the GPS when it sends new coordinates
\item It sends a PUT to the twin HMI on the other server — a \textbf{server-to-server} propagation, not client-to-client
\end{itemize}
\vspace{1cm}
\begin{lstlisting}
updateOtherHMI(): PUT to twin port, payload marked "HMI_UP"
\end{lstlisting}

\end{frame}

\begin{frame}{End-to-End Simulation}
\texttt{SmartAmbulanceSimulation.main} orchestrates a scripted scenario, with \texttt{Thread.sleep} pauses between steps:
\begin{enumerate}
\item Driver sets the temperature in Rear Cabin
\item GPS sends coordinates to the HMI
\item Driver updates emergency/location (propagated to the other HMI)
\item Reliever attempts the same PUT in Cockpit $\rightarrow$ FORBIDDEN
\item Driver switches server and repeats GET/PUT in Cockpit
\item Reliever observes the temperature sensor, then cancels
\end{enumerate}
\end{frame}

% ======================================================
\section{Conclusions}
% ======================================================


\begin{frame}{Questions \& Contacts}
  \vfill
  \begin{center}
    {\Huge Any questions?}

    \vspace{1.5cm}
    {\small E-mail: giuseppe.diblasi@studenti.unipr.it}
  \end{center}
  \vfill
\end{frame}

\end{document}
